\documentclass[11pt]{article}
\usepackage[margin=1in]{geometry}
\usepackage{amsmath,amssymb}
\usepackage{mathtools}
\usepackage{hyperref}
\hypersetup{colorlinks=true, linkcolor=blue, urlcolor=blue}

\newcommand{\rvec}{\mathbf{r}}
\newcommand{\one}{\mathbf{1}}
\newcommand{\wv}{\mathbf{w}'}
\newcommand{\av}{\mathbf{a}'}
\newcommand{\tp}{^{\mathsf T}}
\DeclareMathOperator{\diag}{diag}
\DeclareMathOperator{\rank}{rank}

\title{The over-coordination Hessian and its definiteness}
\author{VBC gate calibration \quad (\texttt{swap\_calibration})}
\date{}

\begin{document}
\maketitle

\section{The energy for $n=v+1$}

We test whether a saturated centre can trap one extra particle. A centre that already
holds its $v$ valency bonds is approached by one more, so there are $n=v+1$ neighbours.
The angles are fixed and the neighbours interact only with the centre, so the only degrees
of freedom are the radial distances $\rvec=(r_1,\dots,r_n)$.

With per-bond counting weight $w(r)$, gate $g(x)$, repulsion $V_r(r)$ and attraction
$V_a(r)$, write the total coordination, the shared gate, the gated single-bond attraction
and its sum as
\[
S=\sum_{j=1}^{n} w(r_j),\qquad
G=g(S),\qquad
a(r)=g\!\left(w(r)\right)\,V_a(r),\qquad
A=\sum_{i=1}^{n} a(r_i).
\]
The total energy is
\[
\boxed{\;E(\rvec)=\sum_{i=1}^{n} V_r(r_i)\;+\;g(S)\sum_{i=1}^{n} a(r_i)\;=\;\sum_i V_r(r_i)+G\,A\;}
\]
The single gate $G=g(S)$ multiplies every bond at once, so adding the extra particle raises
$S$, pulls $G$ down and weakens all bonds together. That collective back-reaction is what
the Hessian must carry.

\medskip
\noindent\textit{Remark.} In the operating range one bond has $w(r)\le 1<n_{\mathrm{cut}}$,
so the single-bond gate saturates, $g(w(r))=1$, and $a(r)=V_a(r)$. We keep $a(r)$ general
in the algebra below and use this simplification only at the end.

\section{Gradient}

Write $g'\equiv g'(S)$, $w'_k\equiv w'(r_k)$, $a'_k\equiv a'(r_k)$. Since
$\partial S/\partial r_k=w'_k$,
\[
\frac{\partial E}{\partial r_k}=V_r'(r_k)+g'(S)\,w'_k\,A+g(S)\,a'_k .
\]
A stationary configuration solves $\partial E/\partial r_k=0$ for every $k$.

\section{The Hessian is diagonal plus rank two}

Differentiating once more (with $g''\equiv g''(S)$) gives, for all $k,l$,
\[
H_{kl}=\frac{\partial^2 E}{\partial r_k\,\partial r_l}
=\delta_{kl}\,D_k
\;+\;g''A\,w'_k w'_l
\;+\;g'\bigl(w'_k a'_l+a'_k w'_l\bigr),
\]
with the diagonal term
\[
D_k=V_r''(r_k)+g'\,w''(r_k)\,A+g\,a''(r_k).
\]
Collect the two coupling vectors $\wv=(w'_1,\dots,w'_n)\tp$ and
$\av=(a'_1,\dots,a'_n)\tp$ into $U=[\,\wv\mid\av\,]\in\mathbb{R}^{n\times2}$ and set
\[
M=\begin{bmatrix} g''A & g' \\[2pt] g' & 0 \end{bmatrix}.
\]
Then the full $n\times n$ Hessian is
\[
\boxed{\;\mathbf{H}=\mathbf{D}+U\,M\,U\tp,\qquad \mathbf{D}=\diag(D_1,\dots,D_n)\;}
\]
Every off-diagonal coupling lives in the rank-two term $U M U\tp$, whose column space is
$\operatorname{span}\{\wv,\av\}$ regardless of how large $n$ is. Written out entrywise,
\[
\mathbf{H}=
\begin{bmatrix}
H_{11} & H_{12} & \cdots & H_{1n}\\
H_{12} & H_{22} & \cdots & H_{2n}\\
\vdots & \vdots & \ddots & \vdots\\
H_{1n} & H_{2n} & \cdots & H_{nn}
\end{bmatrix},
\qquad
\begin{aligned}
H_{kk}&=D_k+g''A\,w'^2_k+2g'\,w'_k a'_k,\\
H_{kl}&=g''A\,w'_k w'_l+g'\bigl(w'_k a'_l+a'_k w'_l\bigr)\quad(k\ne l).
\end{aligned}
\]

\section{Standard definiteness conditions}

The configuration is a locally stable minimum exactly when $\mathbf{H}$ is positive
definite (every eigenvalue positive). It is an energy maximum when $\mathbf{H}$ is negative
definite (every eigenvalue negative), and a saddle, which is also unstable, when
$\mathbf{H}$ is indefinite. So a stable trap means $\mathbf{H}\succ0$, and the arrangement
is unstable whenever $\mathbf{H}\not\succ0$.

\paragraph{(i) Eigenvalue test.}
$\mathbf{H}\succ0$ iff $\lambda_i(\mathbf{H})>0$ for all $i$; $\mathbf{H}\prec0$ iff
$\lambda_i(\mathbf{H})<0$ for all $i$.

\paragraph{(ii) Sylvester's criterion.}
Let $\Delta_k=\det\bigl(\mathbf{H}_{1:k,\,1:k}\bigr)$ be the leading principal minors. Then
\[
\mathbf{H}\succ0 \iff \Delta_k>0\ \ (k=1,\dots,n),
\qquad
\mathbf{H}\prec0 \iff (-1)^k\Delta_k>0\ \ (k=1,\dots,n),
\]
the second being $-\mathbf{H}\succ0$, that is the minors alternate $-,+,-,\dots$. If neither
sign pattern holds the matrix is indefinite.

\paragraph{(iii) Rank-two reduction.}
Because $\rank\bigl(U M U\tp\bigr)\le2$, the $n$-dimensional test collapses. When
$\mathbf{D}\succ0$ (every diagonal term positive), a congruence by $\mathbf{D}^{-1/2}$
followed by the Woodbury identity gives
\[
\mathbf{H}\succ0 \iff M+\bigl(U\tp\mathbf{D}^{-1}U\bigr)^{-1}\succ0,
\]
a single $2\times2$ positive-definiteness test, that is a trace and a determinant. The
$n-2$ directions orthogonal to $\operatorname{span}\{\wv,\av\}$ carry eigenvalue-one factors
and never spoil definiteness once $\mathbf{D}\succ0$. This holds as pure linear algebra
(checked on random $\mathbf{D}\succ0$, agreement $2000/2000$). At a generic radius
$\mathbf{D}$ is usually indefinite, so this shortcut is most useful in the symmetric case
below; otherwise use (i) or (ii).

\section{Symmetric configuration, closed form}

The traps of this energy are symmetric, $r_i=r_\star$ for all $i$. Then $w'_k=w'_\star$,
$a'_k=a'_\star$ and $D_k=D_\star$ are equal, so $\wv=w'_\star\,\one$ and
$\av=a'_\star\,\one$ are both parallel to the all-ones vector $\one$. The rank-two term
becomes rank one and $\mathbf{H}$ diagonalises exactly.
\begin{itemize}
\item \emph{Shape modes}, the subspace $\one^{\perp}$ of dimension $n-1$. Here $UMU\tp$
      annihilates the vector, so
      \[
      \lambda_{\perp}=D_\star,\qquad\text{multiplicity } n-1 .
      \]
\item \emph{Breathing mode}, along $\one$. Using
      $UMU\tp\one=n\bigl(g''A\,w'^2_\star+2g'\,w'_\star a'_\star\bigr)\one$,
      \[
      \lambda_{\parallel}=D_\star+n\,c,\qquad
      c=g''A\,w'^2_\star+2g'\,w'_\star a'_\star .
      \]
\end{itemize}
The whole spectrum is therefore $\{\,D_\star\ (n-1\text{ times}),\ D_\star+nc\,\}$, and the
definiteness reads off at once,
\[
\boxed{\;
\mathbf{H}\succ0\ \text{(stable trap)}\iff D_\star>0\ \text{and}\ D_\star+nc>0,
\;}
\]
\[
\mathbf{H}\prec0\ \text{(maximum, unstable)}\iff D_\star<0\ \text{and}\ D_\star+nc<0,
\qquad\text{indefinite (saddle, unstable) otherwise.}
\]
So the trap is unstable exactly when $D_\star<0$ or $D_\star+nc<0$. The $(n-1)$-fold shape
curvature $D_\star$ is usually the first to change sign, so the stability of the
over-coordinated cluster is governed by the single scalar inequality $D_\star>0$, with
$S=n\,w(r_\star)$, $A=n\,a(r_\star)$ and
\[
D_\star=V_r''(r_\star)+g'(S)\,w''(r_\star)\,A+g(S)\,a''(r_\star)
\;\xrightarrow{\;r_\star>r_{\mathrm{rep}}\;}\;
g'(S)\,w''(r_\star)\,A+g(S)\,a''(r_\star),
\]
the last step using $V_r''=0$ beyond the repulsive core, where a trap sits. The symmetric
radius itself is the single one-dimensional root of the stationarity condition evaluated at
the symmetric point,
\[
V_r'(r_\star)+g'(S)\,w'(r_\star)\,A+g(S)\,a'(r_\star)=0 .
\]
The entire symmetric analysis is thus one 1D root followed by two scalar sign checks, with
no $n$-dimensional optimisation.

\section{Explicit derivatives}

With $\Delta=r_{\mathrm{high}}-r_{\mathrm{low}}$, $s=(r-r_{\mathrm{low}})/\Delta$,
$r_e=(r_{\mathrm{rep}}+r_{\mathrm{att}})/2$, $\sigma=(r_{\mathrm{att}}-r_{\mathrm{rep}})/2$,
$u=(r-r_e)/\sigma$, $z=r/r_{\mathrm{rep}}$ and $t=x-n_{\mathrm{cut}}$, each piece is a
truncated quartic,
\[
\begin{aligned}
w&:\ w=(1-s^2)^2,\ \ w'=\tfrac{-4s(1-s^2)}{\Delta},\ \ w''=\tfrac{4(3s^2-1)}{\Delta^2}
   &&(r_{\mathrm{low}}\le r\le r_{\mathrm{high}}),\\
g&:\ g=(1-t^2)^2,\ \ g'=-4t(1-t^2),\ \ g''=4(3t^2-1)
   &&(n_{\mathrm{cut}}\le x\le n_{\mathrm{cut}}+1),\\
V_a&:\ V_a=E_a(1-u^2)^2,\ \ V_a'=\tfrac{-4E_a u(1-u^2)}{\sigma},\ \ V_a''=\tfrac{4E_a(3u^2-1)}{\sigma^2}
   &&(r_{\mathrm{rep}}< r< r_{\mathrm{att}}),\\
V_r&:\ V_r=E_r(1-z^2)^2,\ \ V_r'=\tfrac{-4E_r z(1-z^2)}{r_{\mathrm{rep}}},\ \ V_r''=\tfrac{4E_r(3z^2-1)}{r_{\mathrm{rep}}^2}
   &&(0\le r\le r_{\mathrm{rep}}),
\end{aligned}
\]
each being constant ($w=1$, $g=1$) or zero outside its window. At the well centre $u=0$,
\[
V_a''(r_e)=\frac{-4E_a}{\sigma^2}=\frac{16\,|E_a|}{(r_{\mathrm{att}}-r_{\mathrm{rep}})^2}\equiv\kappa_0,
\]
the isolated-bond curvature scale (a single decoupled bond reads
$\lambda_{\min}/\kappa_0=1$). For the composite $a(r)=g(w(r))V_a(r)$,
\[
a'=g'(w)\,w'\,V_a+g(w)\,V_a',\qquad
a''=g''(w)\,w'^2V_a+g'(w)\,w''V_a+2g'(w)\,w'V_a'+g(w)\,V_a'',
\]
and in the operating range $w\le1<n_{\mathrm{cut}}$ gives $g(w)=1$, $g'(w)=g''(w)=0$, so
$a=V_a$, $a'=V_a'$, $a''=V_a''$.

\section{Numerical confirmation}

The analytic $\mathbf{H}=\mathbf{D}+U M U\tp$ reproduces the autodiff Hessian
\texttt{jax.hessian(cluster\_energy)} to
$\max|\mathbf{H}_{\text{analytic}}-\mathbf{H}_{\text{autodiff}}|=2.3\times10^{-13}$ over
random configurations. For the looser gate at $v=4$ ($n=5$), the symmetric point
$r_\star=1.31$ gives from the closed form $D_\star=-219.89$ (multiplicity $4$) and
$D_\star+nc=957.05$, matching
$\texttt{eigvalsh}=[-219.89\ (\times4),\ 957.05]$. Here $D_\star<0$, so this radius is a
saddle rather than a stable trap, and the stable trap sits at the nearby radius where
$D_\star$ turns positive.

\section{Eigenvalue localization and gate-specific sign guarantees}

We can say more than the pointwise spectrum. Write the inertia of a symmetric matrix $X$
as $\bigl(n_+(X),n_0(X),n_-(X)\bigr)$, sort its eigenvalues ascending
$\lambda_1(X)\le\cdots\le\lambda_n(X)$, and let $d_1\le\cdots\le d_n$ be the sorted diagonal
curvatures $D_1,\dots,D_n$. The whole analysis rests on one structural fact.

\paragraph{Lemma 1 (the coupling is a two-dimensional saddle).}
\[
\det M=\det\!\begin{bmatrix} g''A & g' \\ g' & 0 \end{bmatrix}=-\,g'(S)^2\le0 .
\]
So whenever $g'(S)\ne0$ the $2\times2$ matrix $M$ has inertia $(1,0,1)$, one positive and
one negative eigenvalue. Provided $\mathbf{w}'$ and $\mathbf{a}'$ are linearly independent,
$UMU\tp$ is therefore a rank-two matrix with exactly one positive eigenvalue, one negative
eigenvalue and $n-2$ zeros. The exceptional parallel case $\mathbf{w}'\parallel\mathbf{a}'$,
which is the symmetric configuration, drops it to rank one and is handled exactly in
Section~5.

\paragraph{Proposition 2 (localization and inertia bounds).}
With $\operatorname{sig}(UMU\tp)=(1,1)$, adding it to $\mathbf{D}$ is a positive rank-one
update followed by a negative rank-one update. Cauchy interlacing applied twice gives, for
every $i$,
\[
d_{i-1}\ \le\ \lambda_i(\mathbf{H})\ \le\ d_{i+1}
\qquad(d_0:=-\infty,\ d_{n+1}:=+\infty),
\]
so each eigenvalue of $\mathbf{H}$ sits within one index of the diagonal curvatures. All but
the smallest and largest are confined to $[d_1,d_n]$, that is at most two eigenvalues leave
the diagonal range, one downward and one upward. In particular
\[
n_-(\mathbf{D})-1\ \le\ n_-(\mathbf{H})\ \le\ n_-(\mathbf{D})+1 ,
\]
and the same for $n_+$. A useful corollary needs no diagonalization,
\[
\mathbf{H}\succ0\ \Longrightarrow\ \text{at most one }D_k<0 ,
\]
so two or more negative diagonal curvatures already force a negative eigenvalue and rule out
a stable trap. Every claim here was checked on $400$ random configurations per gate with no
violation.

\paragraph{Proposition 3 (reduction to one degree of freedom, existence by the intermediate value theorem).}
We do not assume a critical point, we exhibit one. On the symmetric line $r_i=r$ the energy is
a scalar $\varphi(r)=E(r,\dots,r)=n\,V_r(r)+g(nw(r))\,n\,V_a(r)$, and the full $n\times n$
Hessian has at \emph{every} $r$ the exact spectrum $\{D_\star(r)\ (\text{mult }n-1),\
D_\star(r)+nc(r)\}$ of Section~5, with
\[
D_\star(r)=V_r''(r)+g'(S)\,w''(r)\,A+g(S)\,V_a''(r),\qquad S=nw(r),\ A=nV_a(r),
\]
an explicit function of $r$ that uses no stationarity. Two facts follow.
\begin{itemize}
\item $D_\star(r)$ \emph{is not single-signed}. Over the sub-interval where the gate is
      engaged it takes both signs (for $n=5$, looser $[-420,+87]$, optimised $[-775,+42]$), so
      no interval-wide definiteness statement exists. The stability of a symmetric arrangement
      is decided only by the value of $D_\star$ at a critical point, which must first be shown
      to exist.
\item \emph{A symmetric critical point exists}. The reduced gradient
      $\varphi'(r)=n\bigl(g'(S)w'(r)A+g(S)V_a'(r)\bigr)$ is continuous on the gate-engaged
      interval and changes sign there. For $n=5$, looser
      $\varphi'(1.32)/n=-11.80$ and $\varphi'(1.34)/n=+10.03$, optimised
      $\varphi'(1.36)/n=-24.51$ and $\varphi'(1.38)/n=+12.14$, so by the intermediate value
      theorem a root $r_c$ lies in $(1.32,1.34)$ and $(1.36,1.38)$ respectively, exactly one in
      each case.
\end{itemize}

\paragraph{Theorem 4 (classification and the strict separation).}
At the critical point $r_c$ that Proposition 3 guarantees, and only there, stationarity
$g'(S)w'(r_c)A=-g(S)V_a'(r_c)$ collapses the shape curvature to
\[
D_\star(r_c)=g(S)\,\Bigl[V_a''(r_c)-V_a'(r_c)\,\frac{w''(r_c)}{w'(r_c)}\Bigr]=g(S)\,Q(r_c),
\qquad Q(r)\equiv V_a''(r)-V_a'(r)\,\frac{w''(r)}{w'(r)} ,
\]
so $\operatorname{sign}D_\star(r_c)=\operatorname{sign}Q(r_c)$ ($g(S)\ge0$), while the breathing
eigenvalue $D_\star(r_c)+nc(r_c)$ stays positive. The symmetric over-coordinated cluster is thus
never negative definite, it is a stable trap when $Q(r_c)>0$ and an index-$(n-1)$ saddle when
$Q(r_c)<0$. For the base potential $r_{\mathrm{rep}}=1$, $r_{\mathrm{att}}=1.5$, $E_r=20$,
$E_a=-2$, the unique critical point at each valency gives
\[
\begin{array}{c|cccccc}
v=n-1 & 1 & 2 & 3 & 4 & 5 & 6\\\hline
\text{looser } D_\star(r_c)    & -4.18 & 7.21 & 16.01 & 22.60 & 27.79 & 32.09\\
\text{optimised } D_\star(r_c) & -29.66 & -42.98 & -52.90 & -61.09 & -68.21 & -74.59
\end{array}
\]
Hence the \textbf{optimised} gate has $D_\star(r_c)<0$ at every valency, an index-$(n-1)$
saddle, and the extra particle is always rejected, the one positive eigenvalue being the
breathing mode. The \textbf{looser} gate has $D_\star(r_c)>0$ from valency $2$ upward, a genuine
stable trap (valency $1$ is a lone weak index-$1$ saddle). This settles the symmetric family.
Asymmetric critical points are not excluded by the argument, relaxation from scattered initial
placements finds none for these gates, so the symmetric picture appears complete, though that
last step is numerical rather than proved.

\paragraph{Remark (positive-definiteness alone does not separate the gates).}
It would be wrong to say the optimised gate is nowhere positive definite. Both gates admit
positive-definite configurations, on the symmetric line and in the full box. For $n=5$ the
symmetric-line Hessian is positive definite on
$r\in(1.106,1.224)\cup(1.325,1.394)$ for the looser gate and on $r\in(1.367,1.394)$ for the
optimised gate, and a uniform sample of $(1,1.5)^5$ gives comparable positive-definite
fractions, looser about $1.2\%$ and optimised about $1.6\%$, all genuinely bound
($E<0$, gate on). What distinguishes a trap is not the existence of a stable direction but its
coincidence with stationarity. The looser critical point $r_c=1.327$ sits just inside its
positive-definite interval, which begins at $1.325$, whereas the optimised critical point
$r_c=1.366$ sits just below its onset at $1.367$. Both margins in $r$ are small, about $0.0015$,
but $D_\star(r)$ crosses zero steeply, so the shape curvatures there are unambiguous, $+22.6$ for
the looser gate and $-61$ for the optimised gate. A stable trap is stationary and positive
definite at once, and only the looser gate meets both conditions together.

\end{document}
